\subsection{Research Background}

This section situates the study within the broader ESG, ABSA, and multilingual text-analysis landscape, with a focus on environmental, social, and governance (ESG) disclosure, the role of global and Indonesian reporting standards, the structure of sustainability reports and the page-level complexity they introduce for automated analysis, a concise survey of prevailing NLP approaches to ESG analysis and their limitations, and the particular bilingual Indonesian/English challenge in Indonesian corporate reporting. The synthesis integrates the cited frameworks and methodological considerations to motivate a robust ABSA-based investigation of Indonesian disclosures.

\subsubsection{ESG framework definitions and GRI/SASB standards}
ESG as a framework for risk, value creation, and accountability. ESG measures provide a lens through which firms’ environmental, social, and governance activities are assessed for risk management, strategic planning, and investor decision-making. A body of literature and practice emphasizes that ESG disclosures aim to illuminate how firms create or erode long-term value while signaling governance quality and risk management to capital markets \cite{10.1177/1086026620945342,10.1111/jacf.12411,10.1016/j.cpa.2021.102309}. The underlying rationale is that standardized ESG information improves transparency, comparability, and market efficiency, thereby supporting responsible investing and regulatory compliance \cite{10.1177/1086026620945342,10.20900/jsr20210006,10.20944/preprints202407.2036.v1}

Core reporting standards and their theoretical alignment. Globally, a dominant tension exists among major reporting frameworks—GRI (global stakeholder-oriented), SASB (industry-specific, financially material focus), TCFD (climate-risk centric), CDSB, CDP, and the emerging ISSB standards under IFRS. These bodies have pursued alignment through initiatives like the Corporate Reporting Dialogue and joint guidance to map SASB, GRI, and CDP against TCFD, with a continuing push toward a unified global standard set while acknowledging residual fragmentation \cite{10.1108/medar-11-2020-1097,10.20900/jsr20210006,10.1108/aaaj-06-2021-5330,10.1108/aaaj-02-2020-4445,10.1016/j.cpa.2021.102309,10.1111/1911-3838.12232}. The literature notes that SASB emphasizes financially material disclosures for investors, whereas GRI adopts a broader stakeholder perspective, highlighting the materiality tension across frameworks and the need for harmonization to reduce reporting overload \cite{10.20900/jsr20210006,10.1016/j.cpa.2021.102309}

Indonesian regulatory context (global and local). Regulatory convergence toward formalized ESG disclosure in Indonesia is framed by global expectations and local authorities. Indonesia’s market uses a mix of standards and regulatory expectations—OJK (the Financial Services Authority) rules and IDX (the Exchange) disclosure requirements—creating a structured environment for ESG information that is publicly accessible and regulated. The regulatory regime encourages formalized, traceable disclosures across filings and annual reports, signaling a demand for robust analytical tooling to interpret this information coherently for market participants \cite{10.22617/tcs210472-2,10.1108/mf-01-2023-0020,10.1007/s40804-021-00233-z} . This regime aligns with global moves toward standardized ESG reporting but also elevates the complexity of local lexicons and regulatory mappings.

Implications for ABSA. The regulatory density and standardization ambitions imply that ESG disclosures are not merely narrative; they embed regulatory motifs, risk signals, and governance cues that require granular, aspect-level interpretation. ABSA—which links sentiment to specific ESG aspects—offers a mechanism to translate regulatory language into granular, decision-useful signals (e.g., which environmental topics or governance mechanisms drive risk signaling or capital-market responses) \cite{10.1002/csr.70089,10.69593/ajbais.v4i04.130,10.1108/mf-01-2023-0020,10.1007/s40804-021-00233-z,10.1111/eufm.12509}

How sustainability reports are structured and the page-level complexity this creates for automated analysis
Typical sustainability-report structure. Sustainability reports commonly adopt a multi-section layout aligned with material topics, governance disclosures, strategy and risk management, and performance metrics, often organized by ESG pillar and then by sub-topics (e.g., emissions, energy efficiency, supply chain governance, diversity and inclusion, board oversight). This structure supports a narrative of materiality, strategy alignment, and accountability, yet it also yields dispersed content across sections, tables, figures, and appendices, complicating automated extraction and cross-document synthesis \cite{10.1111/jacf.12411,10.1111/1911-3838.12232,10.1111/ablj.12210} .

Cross-framework implications for document structure. When firms report under multiple frameworks (GRI, SASB, TCFD, etc.), the same information may be presented in different templates, with varying indicator identifiers, materiality thresholds, and regulatory lexicons. The result is a heterogeneous document design in which regulatory cues, indicator mappings, and narrative themes appear across sections and even within the same sentence. This heterogeneity challenges parsers that rely on consistent headings or standardized ontologies, and it motivates design choices that can align textual content to a unified set of ESG topics and regulatory motifs \cite{10.1108/medar-11-2020-1097,10.20900/jsr20210006,10.31436/iiumlj.v30is2.755,10.1787/a9889ba3-en}

Page-level complexity for ABSA pipelines. For ABSA, the page-level complexity manifests in several ways:

Multi-topic co-location: Sentiment toward an environmental topic may appear in close proximity to social or governance content, requiring robust topic disambiguation.
Lexical variation and terminology drift: Terminology shifts across frameworks and languages necessitate adaptable lexicons and context-sensitive models.
Table and figure content: Quantitative indicators in tables/figures often accompany qualitative narratives; extracting numerical signals alongside textual sentiment requires multi-modal integration.
Regulatory mapping: Cross-referencing narrative claims with regulatory requirements (e.g., climate risk disclosures, governance practices) demands models that can align textual cues to regulatory lexicons and standards.
Language and translation effects: In Indonesia, reports frequently include Indonesian and English content within the same document or across companion reports, complicating alignment and requiring cross-language capabilities (see Section 4).
Implications for ABSA design. To bridge the depth and alignment gaps, ABSA pipelines must (i) operate on multilingual corpora, (ii) incorporate topic- and framework-aware lexicons, (iii) support aspect-term extraction and polarity scoring at fine granularity, and (iv) provide explainable outputs that map sentiments to regulatory themes and material ESG topics. The literature indicates that ABSA, when combined with cross-lingual NLP and explainability techniques, can yield interpretable, action-oriented signals from complex ESG disclosures \cite{10.3390/math12193119,10.1002/csr.70089,10.1111/eufm.12509,10.1109/access.2023.3348342,10.2139/ssrn.3738629}

\subsubsection{Brief survey of existing NLP approaches to ESG analysis and their limitations}
Traditional sentiment and topic analysis. Early ESG NLP work emphasized document-level sentiment, topic modeling (e.g., LDA), and keyword extraction to summarize disclosures. While these methods illuminate broad sentiment trends and dominant topics, they lack the granularity needed for investor due diligence and regulatory monitoring, as they do not tie sentiment to specific ESG aspects or regulatory criteria \cite{10.3390/su15054134,10.69593/ajbais.v4i04.130,10.1111/eufm.12509}

Aspect-based sentiment analysis (ABSA). ABSA addresses the granularity gap by extracting aspect terms, associating sentiment polarity with those aspects, and yielding a structured representation of ESG narratives. ABSA enables more precise risk/opportunity signaling and supports explainability by linking opinions to concrete topics (e.g., emissions, governance processes). Recent systematic reviews highlight ABSA’s strength in extracting fine-grained signals but also note challenges in domain adaptation, cross-domain labeling, and interpretability in highly technical ESG texts \cite{10.1002/csr.70089,10.3390/math12193119,10.2139/ssrn.3738629,10.1111/eufm.12509,10.1109/access.2023.3348342}

Cross-lingual and multilingual NLP in ESG. Given bilingual or multilingual reporting (e.g., Indonesian/English), cross-lingual ABSA methods leverage translation, multilingual pre-trained models (e.g., mBERT, XLM-R), and alignment-based labeling to enable cross-language sentiment analysis. Surveys indicate cross-lingual ABSA effectiveness when annotated multilingual data exist or when zero-shot transfer is leveraged; however, performance gaps persist when domain-specific terminology or regulatory lexicons are sparse in high-resource languages \cite{10.1016/j.inffus.2025.103073,10.1111/eufm.12509,10.48550/arxiv.2511.00024}

Explainability and regulatory relevance. A recurring limitation across NLP methods is model opacity. To satisfy regulatory and investor demands for transparency, researchers increasingly integrate explainable AI (XAI) techniques (e.g., SHAP, LIME) with ABSA outputs to justify aspect-level sentiment decisions and highlight salient features driving the judgments. Systematic reviews of ESG analytics emphasize the need for interpretability, bias mitigation, and robust evaluation when deploying NLP in financial contexts \cite{10.1002/csr.70089,10.2139/ssrn.3738629,10.1108/mf-01-2023-0020}

Indonesian-specific and regional gaps. While global ABSA and multilingual NLP research is burgeoning, the Indonesian disclosure landscape presents domain-specific lexicons, industry sectors, and regulatory language that are underrepresented in high-quality multilingual training data. This creates a critical gap that a bilingual ABSA framework must address through domain-adaptive fine-tuning, lexicon customization, and regulatory lexicon integration to ensure reliable extraction of ESG topics and sentiment from Indonesian corporate reporting \cite{10.69593/ajbais.v4i04.130,10.1108/mf-01-2023-0020,10.1111/eufm.12509}

Summary of limitations. Collectively, the literature identifies (i) a need for fine-grained, aspect-level signals linked to financially material ESG topics; (ii) challenges in cross-language interpretation and translation-induced variability; (iii) the demand for explainable outputs aligned with regulatory criteria; and (iv) data standardization and domain adaptation issues that hinder cross-framework comparability. These gaps motivate the current focus on ABSA-based, multilingual ESG analysis with an emphasis on Indonesian disclosures, regulatory alignment, and interpretability \cite{10.1002/csr.70089,10.2139/ssrn.3738629,10.1108/mf-01-2023-0020,10.1016/j.inffus.2025.103073,10.1111/eufm.12509,10.48550/arxiv.2511.00024,10.3390/math12193119}

\subsubsection{The specific bilingual Indonesian/English challenge in Indonesian corporate reporting}
Prevalence of bilingual reporting. Indonesian corporate reporting frequently features content in both Indonesian and English, either within sustainability reports or across companion documents (e.g., annual reports and separate ESG disclosures). This bilingual presentation creates cross-language alignment demands, complicating standard NLP pipelines that assume monolingual input and stable terminologies. Cross-lingual ABSA research specifically focuses on scenarios where multilingual content necessitates either robust multilingual models or effective translation-then-analyze strategies to preserve nuance across languages \cite{10.1016/j.inffus.2025.103073,10.1111/eufm.12509,10.48550/arxiv.2511.00024}

Language phenomena and regulatory lexicons. Indonesian-English reporting surfaces code-switching, mixed-language terminology, and varying regulatory lexicons that reflect both local compliance language and international standard terminology. Cross-language ABSA must handle coinages, acronyms, and sector-specific terms that differ in Indonesian and English, requiring dynamic lexicon expansion and contextual embedding strategies that can capture register and domain semantics across languages \cite{10.1016/j.inffus.2025.103073,10.1111/eufm.12509,10.48550/arxiv.2511.00024}

Model design implications. For bilingual Indonesian/English ESG ABSA, model design choices include (i) employing multilingual pre-trained transformers with domain-adaptive fine-tuning, (ii) integrating translation or alignment components to unify cross-language sentiment signals, (iii) constructing regulatory-anchored aspect taxonomies that map Indonesian and English ESG terms to common topics and to the regulatory frameworks (GRI, SASB, TCFD), and (iv) ensuring explainability to satisfy both regulators and investors. Empirical evidence from cross-lingual ABSA surveys indicates that superior performance arises from multilingual models with cross-language alignment and explicit aspect annotation, rather than naïve translation-only approaches \cite{10.1016/j.inffus.2025.103073,10.1111/eufm.12509,10.48550/arxiv.2511.00024,10.3390/math12193119}

Practical relevance for Indonesian markets. Indonesian ABSA that can robustly parse bilingual reports supports regulatory monitoring (e.g., alignment with OJK/IDX expectations), investor due diligence, and risk management by enabling granular extraction of sentiment signals tied to specific ESG aspects. The literature already stresses the importance of regulatory-aligned, granular ESG analytics for governance and market outcomes, while highlighting the need for multilingual, explainable approaches in non-English reporting contexts \cite{10.22617/tcs210472-2,10.1108/mf-01-2023-0020,10.1007/s40804-021-00233-z,10.1002/csr.70089,10.69593/ajbais.v4i04.130,10.1111/eufm.12509,10.1109/access.2023.3348342}

Transition to conceptual and methodological motivation
Why ABSA is particularly well suited to Indonesian ESG disclosures. ABSA provides a structured representation of ESG narratives at the aspect level, enabling investors and regulators to detect nuanced risk signals, governance weaknesses, or positive commitments that are otherwise obscured in document-level summaries or coarse scoring. In multilingual Indonesian/English contexts, ABSA is complemented by cross-lingual NLP methods and multilingual transformers to unify cross-language sentiment signals around shared environmental, social, and governance topics, which is essential for comparability across domestic and international stakeholders \cite{10.1002/csr.70089,10.1016/j.inffus.2025.103073,10.1111/eufm.12509,10.48550/arxiv.2511.00024,10.3390/math12193119}

Addressing the page-level complexity with ABSA. By anchoring sentiment to explicit ESG aspects and linking those aspects to regulatory themes, ABSA can reduce the information extraction burden from dense documents. It supports explainable outputs that tie sentiment to concrete sections, indicators, and standards, responding to the demand for transparency in ESG analytics described in the literature and discussed in the Indonesian regulatory milieu \cite{10.1002/csr.70089,10.2139/ssrn.3738629,10.1108/mf-01-2023-0020,10.1108/medar-11-2020-1097}

Gaps this research aims to close. Although substantial work demonstrates the potential of ABSA and cross-lingual NLP for ESG analysis, there remains a scarcity of rigorously evaluated, bilingual ABSA pipelines tailored to Indonesian corporate reporting, with explicit alignment to local regulatory lexicons and international standards. This study seeks to fill that gap by developing an ABSA framework capable of (i) handling Indonesian-English bilingual texts, (ii) mapping sentiments to GRI/SASB/TCFD-aligned aspects, and (iii) providing explainable, regulator-friendly outputs suitable for Indonesian stakeholders and global investors.

\subsubsection{Summary}
The ESG landscape is governed by converging but heterogeneous standards (GRI, SASB, TCFD, ISSB, and others), creating a complex, multi-framework disclosure ecosystem that Indonesian firms navigate under OJK/IDX oversight. This motivates the need for fine-grained ABSA that can interpret disclosures at the aspect level and map sentiments to regulatory themes, thereby supporting better risk assessment, governance monitoring, and investment decision-making \cite{10.1108/medar-11-2020-1097,10.20900/jsr20210006,10.1108/aaaj-06-2021-5330,10.1108/aaaj-02-2020-4445}. The structural characteristics of sustainability reports—multi-section, cross-framework content, and mixed Indonesian/English language—imply page-level complexity that challenges traditional NLP, reinforcing the case for ABSA with explainability and cross-lingual capabilities \cite{10.1111/ablj.12210,10.1111/jacf.12411,10.1016/j.cpa.2021.102309,10.1111/1911-3838.12232}. Existing NLP approaches offer promising directions but exhibit limitations in granularity, cross-language adaptation, and regulatory alignment, particularly in Indonesian contexts \cite{10.1002/csr.70089,10.69593/ajbais.v4i04.130,10.1111/eufm.12509,10.3390/math12193119}. The bilingual Indonesian/English reporting environment further intensifies these challenges, necessitating a bilingual ABSA approach informed by cross-lingual transfer, lexical adaptation, and regulatory mapping \cite{10.1016/j.inffus.2025.103073,10.1111/eufm.12509,10.48550/arxiv.2511.00024}

The transition to the research core is thus warranted: to design, implement, and evaluate a bilingual ESG ABSA pipeline tailored to Indonesian disclosures, capable of delivering aspect-level sentiment signals aligned with GRI/SASB/TCFD concepts, while providing explainable outputs suitable for regulators and investors. The ensuing chapters (Research Design and Data, Methodology, and Experiments) will detail the data sources, annotation schema, model architectures, evaluation metrics, and case studies that operationalize these objectives.

Annotated bibliography (machine-readable XML)